\documentclass[11pt]{article}
\usepackage[T1]{fontenc}
\usepackage{lmodern,amsmath,amssymb,amsthm,geometry,microtype}
\geometry{margin=1in}
\usepackage[colorlinks=true,urlcolor=blue,citecolor=blue,linkcolor=blue]{hyperref}
\newtheorem{theorem}{Theorem}
\newtheorem{lemma}[theorem]{Lemma}
\newtheorem{proposition}[theorem]{Proposition}
\newcommand{\Pone}{\mathbb P^1}
\newcommand{\D}{\mathbb D}
\newcommand{\ddc}{dd^c}
\newcommand{\FS}{\omega_{\mathrm{FS}}}
\newcommand{\ind}{\mathbf 1}
\title{An interior second-order defect from a Hele--Shaw contact arc}
\author{Alec Kriebel\\\small Independent researcher\\
\small\href{https://orcid.org/0009-0001-9320-500X}{ORCID: 0009-0001-9320-500X}}
\date{15 October 2026}
\begin{document}
\maketitle
\begin{abstract}
We give a detailed application of Ross--Witt Nystr\"om's smooth Hele--Shaw
contact-arc construction: the associated homogeneous complex Monge--Amp\`ere
solution on $\Pone\times\D$ has no second-order real Taylor expansion at an
interior point, despite smooth strictly K\"ahler fibre boundary data. Filling
a zero-area contact arc gives a strict Green-function time-slope gap; smooth
variation on the covering surface supplies finite nonzero curvature on the
other side. Legendre duality then yields incompatible quadratic coefficients.
The 2021 Oberwolfach report already announces an interior construction.
This is an attributed proof/application note; historical priority of the
deduction is unresolved and no first-discovery claim is made.
\end{abstract}

\section{Statement and attribution}
Let $\D=\{\tau\in\mathbb C:|\tau|<1\}$ and let $\pi:\Pone\times\D\to\Pone$
be projection. We use the coherent normalization
\[
 d^c=\frac{i}{4\pi}(\bar\partial-\partial),\qquad
 \ddc=\frac{i}{2\pi}\partial\bar\partial,\qquad
 \ddc\log|z|^2=\delta_0,\qquad
 \FS=\ddc\log(1+|z|^2),\quad \int_{\Pone}\FS=1.
\]
A second-order real Taylor expansion at $p$ means an expansion in real local
coordinates of the form $u(p+h)=u(p)+\ell(h)+Q(h)+o(|h|^2)$, where $\ell$ is
linear and $Q$ is quadratic.

\begin{theorem}\label{thm:main}
There exists $\phi\in C^\infty(\Pone)$ with
$\omega_\phi:=\FS+\ddc\phi>0$ such that the bounded weak solution $\Phi$ of
\begin{equation}\label{eq:pde}
 \pi^*\FS+\ddc\Phi\geq0,\qquad
 (\pi^*\FS+\ddc\Phi)^2=0,\qquad
 \Phi(z,\tau)=\phi(\tau z)\quad (|\tau|=1)
\end{equation}
has no second-order real Taylor expansion at a point $(z_*,\tau_*)$ with
$0<|\tau_*|<1$. The equation is in the Bedford--Taylor sense, and each
boundary fibre form is strictly positive. In particular $\Phi$ is not $C^2$
in the interior.
\end{theorem}

The construction of $\phi$, the obstacle-current identity, the smooth strong
flow on a covering surface, and the Monge--Amp\`ere/Hele--Shaw duality are
all due to Ross--Witt Nystr\"om \cite{applications,harmonic}.
The additional argument written here is the arc-filled Green comparison,
its combination with finite left curvature, and the resulting interior
Taylor defect. Green extremality, the Hopf lemma and convex conjugacy used
in that argument are classical tools.

This attribution is material to the status of the result. Their survey
\cite[\S4.3]{survey} discusses the interior-location question in its dated
2017/2018 context. Witt Nystr\"om's joint-work report with Ross
\cite[pp.~1323--1324]{owr} already announces and sketches an interior
construction, using a terminal-domain harmonic-disc route and an inverse-flow
density whose smoothness is left tentative there; that discussion permits
vanishing density at the terminal singular points. It must receive credit
for the announced phenomenon. The present note uses instead a finite-time
contact arc and the strictly positive form supplied by
\cite[Proposition~3.4]{applications}. The earlier published contact theorem
\cite[Theorem~6.2]{applications} states its nonsmoothness conclusion on the
boundary circle. The rank-degeneration examples in \cite{rank} by themselves
do not give the Taylor defect below.

A bounded literature audit through 15 October 2026 found no complete exact
predecessor in the recorded inspected scope. It does not establish historical
firstness, independent discovery, or present worldwide openness. In particular,
the Oberwolfach contribution is not being recast as a numbered unresolved
conjecture. This note supplies a checkable proof application, rather than a
claim to a newly discovered general phenomenon. Source versions, precise
inspection scope and remaining coverage/access gaps are in the supplement.

\section{Credited flow and duality inputs}
For a potential $\phi$ as above, set
\[
 \psi_t=\sup\{\psi\in\operatorname{PSH}(\Pone,\FS):
            \psi\leq\phi,\ \nu_0(\psi)\geq t\},\qquad
 \Omega_t=\{\psi_t<\phi\},\quad 0\leq t\leq1.
\]
Here $\nu_0$ is the coefficient of the unit logarithmic pole, with the usual
upper-semicontinuous envelope convention. The domains increase with $t$.
The envelopes decrease and are concave in $t$. For $0<t<1$,
\begin{equation}\label{eq:current}
 \FS+\ddc\psi_t=\ind_{\Omega_t^c}\omega_\phi+t\delta_0,
 \qquad \int_{\Omega_t}\omega_\phi=t,
\end{equation}
and $\psi_t$ is $C^{1,1}$ away from the pole
\cite[Proposition~2.5]{applications}.

We use the following specific features of the published construction,
not an assertion about arbitrary self-contacting flows.

\begin{proposition}[Ross--Witt Nystr\"om construction]\label{prop:input}
One can choose a smooth strictly positive $\omega_\phi$ and $T\in(0,1)$
so that $\Omega=\Omega_T$ is connected and its boundary has a smooth
contact segment $S$ of positive length. At a relative interior point
$q\in S$, a sufficiently small coordinate disc $B$ satisfies
$B\setminus S\subset\Omega$, with fluid on both sides of $S$.
The pre-contact domains are projections of a smooth family of smoothly
bounded simply connected strong-flow domains $\Omega'_t$ on a fixed covering
surface, extending smoothly through $t=T$. Projection is conformal and
injective on the open domains for $t\leq T$ close to $T$.
The limiting outward boundary velocity on the lift is strictly positive.
\end{proposition}

These are the construction in \cite[Proposition~3.4 and its proof]{applications},
with the strong-flow extension from its Theorem~2.15 (and the two-sided local
flow discussed in \cite[Remark~7.1]{moduli}). Nonvanishing limiting velocity
and the two-sided fluid geometry are recorded in the proof of
\cite[Theorem~6.2]{applications}. The covering resolves the contact:
distinct boundary points can project to the same arc. Injectivity is required
only on the open fluid domain, not on its closure. A smooth compact domain
and collar on the cover can be placed in a conformal coordinate chart;
smooth Green dependence is then the uniformly elliptic Dirichlet dependence
of \cite[Appendix~A]{applications}.

Write $G_U(w,0)$ for the negative Green function with singular part
$\log|w|^2$ at the injection point. The duality
\cite[Theorem~4.4]{applications} gives
\begin{equation}\label{eq:dual}
 \widetilde\Phi(w,\tau)
 =\sup_{0\leq t\leq1}\{\psi_t(w)+(1-t)\log|\tau|^2\}.
\end{equation}
For $\tau\ne0$, \cite[Proposition~4.2]{applications} identifies it with
the bounded solution in \eqref{eq:pde} through
\begin{equation}\label{eq:twist}
 \Phi(z,\tau)+\log|\tau|^2+\log(1+|z|^2)
 =\widetilde\Phi(\tau z,\tau)+\log(1+|\tau z|^2).
\end{equation}
We use this identity in finite affine coordinates near the point in question.
The solution is the usual envelope with the prescribed smooth boundary data.
A constant $\min\phi$ is an admissible lower bound; vertical subharmonic
comparison gives the upper bound $\max\phi$. The cited duality concerns this
weak solution, with its boundary trace, not just a formal transform.

\section{A strict Green-function slope gap}
Choose $q,B$ as in Proposition~\ref{prop:input}, away from $0$ and the arc
endpoints, and put $D_+=\Omega\cup B$. This is an open connected proper
domain. Its difference from $\Omega$ has area zero, whereas its complement
has positive $\omega_\phi$-area $1-T$. Thus both domains are Greenian.

\begin{lemma}\label{lem:gap}
For every $w\in\Omega\setminus\{0\}$,
\[
 \partial_t^+\psi_T(w)\leq G_{D_+}(w,0)<G_\Omega(w,0).
\]
\end{lemma}
\begin{proof}
On $D_+\setminus\{0\}$, \eqref{eq:current} makes $\psi_T$ plus a local
potential of $\FS$ distributionally harmonic: the newly filled part of $S$
has area zero, and the current identity supplies the entire measure, including
the absence of a line measure. Weyl regularity therefore gives harmonic
continuation across the arc. This is not an arbitrary removability assertion
for positive-capacity sets.

For $0<h<1-T$ let $v_h=(\psi_{T+h}-\psi_T)/h$. Then $v_h\leq0$ and,
on $D_+$,
\[
 \ddc v_h=\delta_0+h^{-1}\ind_{\Omega_{T+h}^c}\omega_\phi\geq\delta_0.
\]
In fact monotonicity of the domains makes the extra term zero almost
everywhere on $D_+$; the weaker inequality suffices. The logarithmic
coefficient at $0$ is one. The extremal characterization of the negative
Green function gives $v_h\leq G_{D_+}$.

Domain monotonicity gives $G_{D_+}\leq G_\Omega$. The inequality is strict:
the old Green function tends to zero at a smooth interior point of the
contact slit from either side, as is also seen on the smooth lift.
Meanwhile $G_{D_+}(q,0)<0$ since $q$ is now an interior point. Their
nonnegative harmonic difference on $\Omega$, with the equal logarithmic
poles cancelled, is nonzero. The strong maximum principle makes it positive
everywhere. Finally $t\mapsto\psi_t(w)$ is finite and concave near $T$, so
the right derivative exists and passing $h\downarrow0$ proves the claim.
\end{proof}

The arc must be nonpolar. Filling a single isolated point would not justify
this strict comparison. No description, smoothness or filling-rate estimate
for the post-contact weak flow was used.

\section{Finite nonzero curvature from the smooth lift}
Fix a finite affine point $w_*\in\Omega\setminus\{0\}$ and write
$a(t)=\psi_t(w_*)$. It belongs to the smoothly evolving domains for all
$t<T$ sufficiently close to $T$.

\begin{lemma}\label{lem:left}
The pre-contact branch satisfies
\[
 a'_-(T)=\kappa_0:=G_\Omega(w_*,0)<0,\qquad
 a''_-(T)=L\in(-\infty,0),
\]
with smooth one-sided dependence up to $T$.
\end{lemma}
\begin{proof}
On a smooth pre-contact domain, differentiating the harmonic Dirichlet
problem for $\psi_t$ gives
\[
 \partial_t\psi_t(w)=G_{\Omega_t}(w,0).
\]
Indeed this derivative is harmonic away from its fixed unit logarithmic
pole. Its boundary value is zero: differentiating $\psi_t-\phi=0$ along
the moving boundary has no spatial contribution, by the $C^1$ smooth fit.
Uniqueness of the Green Dirichlet problem proves the identity.
The smooth Dirichlet problem on the lifted family justifies these
differentiations up to the one-sided limit at $T$.

Choose fixed lifts of $0$ and $w_*$ under the fixed covering projection.
For $t\leq T$ the lifted and original Green functions agree by conformal
invariance. Smooth domain dependence on the lift gives a finite derivative
$H=\partial_tG_{\Omega'_t}|_{T}$, harmonic even at the pole after removal
of its time-independent logarithmic singularity. Differentiating its zero
boundary value gives
\[
 H=-V\partial_nG_{\Omega'_T}\quad\hbox{on }\partial\Omega'_T.
\]
Here $V>0$ is the limiting outward normal velocity. The outward normal
derivative of the negative Green function is positive by Hopf's lemma.
Thus $H<0$ on the boundary and throughout the domain, and evaluation at
the fixed lift of $w_*$ gives finite $L<0$. The projected singular boundary
is never used as a smooth Dirichlet boundary at $T$.
\end{proof}

\section{Legendre plateau and incompatible Taylor coefficients}
Lemma~\ref{lem:gap} supplies
\begin{equation}\label{eq:slopes}
 a'_+(T)\leq b:=G_{D_+}(w_*,0)<\kappa_0=a'_-(T).
\end{equation}
For $\kappa$ near $\kappa_0$ set
\[
 F(\kappa)=\widetilde\Phi(w_*,e^{\kappa/2})
 =\kappa+\sup_{0\leq t\leq1}\{a(t)-\kappa t\}.
\]
Global concavity is essential. For $b<\kappa<\kappa_0$, the supporting
slopes in \eqref{eq:slopes} make $T$ the unique global maximizer, so
\begin{equation}\label{eq:affine}
 F(\kappa)=a(T)+(1-T)\kappa.
\end{equation}
For $\kappa>\kappa_0$ sufficiently close, Lemma~\ref{lem:left} and the
implicit function theorem give a unique $t(\kappa)<T$ near $T$ with
$a'(t(\kappa))=\kappa$. It too is the unique global maximizer: a supporting
line for a concave function is global, and local strict curvature excludes
any second point of equality. Thus no unknown post-contact branch can
overtake it. Endpoint values in the supremum cause no difficulty, even
if interpreted in the extended-real envelope convention.

Put $\delta=\kappa-\kappa_0$ and $x=t(\kappa)-T$. As $\delta\downarrow0$,
\[
 x=\delta/L+o(\delta),\qquad
 a(T+x)=a(T)+\kappa_0x+\tfrac12Lx^2+o(x^2).
\]
Substitution, together with \eqref{eq:affine} on the other side, yields
\begin{equation}\label{eq:quadratic}
 F(\kappa_0+\delta)-F(\kappa_0)-(1-T)\delta
 =\begin{cases}
 0,&\delta<0\text{ sufficiently small},\\[2pt]
 -\delta^2/(2L)+o(\delta^2),&\delta>0.
 \end{cases}
\end{equation}
Since $L<0$, no common quadratic Taylor coefficient exists. This proves
failure of a second-order expansion, rather than merely a discontinuous
limiting Hessian.

\begin{proof}[Completion of Theorem~\ref{thm:main}]
The number $\kappa_0$ is finite and negative. Define
$\tau_*=e^{\kappa_0/2}\in(0,1)$ and $z_*=w_*/\tau_*$.
Along the smooth curve $(z,\tau)=(w_*e^{-\kappa/2},e^{\kappa/2})$,
\eqref{eq:twist} gives
\[
 \Phi(w_*e^{-\kappa/2},e^{\kappa/2})
 =F(\kappa)+\log\frac{1+|w_*|^2}{e^\kappa+|w_*|^2}.
\]
The added term is smooth. A second-order real Taylor expansion for $\Phi$
at $(z_*,\tau_*)$ would restrict to one along this curve, contradicting
\eqref{eq:quadratic}. Finally $(z,\tau)\mapsto\phi(\tau z)$ is jointly
smooth on $\Pone\times\partial\D$, and circle rotations preserve $\FS$,
so each boundary fibre form is strictly positive.
\end{proof}

\paragraph{Scope and verification.}
This is an existence recipe using a credited designer-flow theorem,
not a closed-form potential. The argument requires a positive-length arc,
$0<T<1$, a non-pole interior observation point, the smooth lifted family,
and finite strictly negative $L$. It does not prove an arbitrary point
tangency or terminal-time slit has this defect. The supplementary exact
rational verifier checks the scalar support and quadratic signs, including
invalid-hypothesis controls; it does not certify the geometric PDE inputs
by computation. The proof above supplies those analytic deductions from
the cited results.

\paragraph{AI use and review status.}
AI tools were used extensively in investigating the problem, deriving,
drafting and checking the argument, auditing sources and preparing the
verification materials. Independent AI adversarial reviews are recorded
with their actual scope in the repository. This preprint is unrefereed and
has not undergone conventional human peer review. The author is responsible
for the claims and their qualifications.

\begin{thebibliography}{9}
\bibitem{applications} J. Ross and D. Witt Nystr\"om,
\emph{Applications of the duality between the homogeneous complex
Monge--Amp\`ere equation and the Hele--Shaw flow},
Ann. Inst. Fourier \textbf{69} (2019), 1--30.
\href{https://doi.org/10.5802/aif.3237}{doi:10.5802/aif.3237}.
\bibitem{harmonic} J. Ross and D. Witt Nystr\"om,
\emph{Harmonic discs of solutions to the complex homogeneous
Monge--Amp\`ere equation}, Publ. Math. IHES \textbf{122} (2015), 315--335.
\href{https://doi.org/10.1007/s10240-015-0074-0}{doi:10.1007/s10240-015-0074-0}.
\bibitem{moduli} J. Ross and D. Witt Nystr\"om,
\emph{The Hele--Shaw flow and moduli of holomorphic discs},
Compos. Math. \textbf{151} (2015), 2301--2328.
\href{https://doi.org/10.1112/S0010437X15007526}{doi:10.1112/S0010437X15007526}.
\bibitem{rank} J. Ross and D. Witt Nystr\"om,
\emph{On the maximal rank problem for the complex homogeneous
Monge--Amp\`ere equation}, Anal. PDE \textbf{12} (2019), 493--503.
\href{https://doi.org/10.2140/apde.2019.12.493}{doi:10.2140/apde.2019.12.493}.
\bibitem{survey} J. Ross and D. Witt Nystr\"om,
\emph{The Dirichlet problem for the complex homogeneous Monge--Amp\`ere
equation}, author version
\href{https://arxiv.org/abs/1712.00405v2}{arXiv:1712.00405v2} (2018).
\bibitem{owr} D. Witt Nystr\"om (joint work with J. Ross),
\emph{Singular solutions to the homogeneous Monge--Amp\`ere equation via
the Hele--Shaw flow}, in \emph{Geometric Methods of Complex Analysis},
Oberwolfach Report 25/2021, pp.~1321--1325.
\href{https://doi.org/10.4171/owr/2021/25}{doi:10.4171/owr/2021/25}.
\end{thebibliography}
\end{document}
